\section{Proofs of the inherited NBO results}\label{supp:proofs67}
These proofs retain the original statements and arguments. Their placement in the supplement changes the reading order, not their hypotheses.

\subsection{Proof of the result \ref{prop:sandwich49}}\label{supp:proof-core49-1}
\begin{proof}
Backward comparison and cash invariance give
\[
 V_t^*\ge f_t-\sum_{k=t}^{T-1}B_{t,k}u_k-B_{t,T}u_T,
 \quad
 J_t^\pi\le f_t+\sum_{k=t}^{T-1}B_{t,k}v_k-B_{t,T}\ell_T,
\]
where $B_{t,k}=\prod_{j=t}^{k-1}\beta_j$. The terminal claim is the stipulated enclosure; applying the Bellman or fixed-policy operator proves each inductive step. Subtraction gives \eqref{eq:sandwich49}. A shift $f_t\mapsto f_t+c_t^0$ adds $c_t^0-\beta_tc_{t+1}^0$ to $u_t$ and subtracts it from $v_t$. The terminal width is unchanged.
\end{proof}

\subsection{Proof of the result \ref{thm:construct49}}\label{supp:proof-core49-2}
\begin{proof}
The selected node label lies between the node Bellman infimum minus $e_t$ and that infimum plus $D_tk_t+e_t$. The Lipschitz property of the infimum implies that every cone lies above it minus $e_t$. At a nearest node its cone exceeds the infimum by at most $D_tk_t+e_t+2L_th_t$. At the attaining original index, the two feasible pairs and \eqref{eq:repair49} give
\[
 Q_t(f_{t+1};x,\Gamma_{t,i}(x))
 \le y_{t,i}+e_t+L_t\|x-x_{t,i}\|_1+D_t\zeta_t.
\]
This proves \eqref{eq:enclosures49}. The terminal residual belongs to $[-e_T,e_T+2L_gh_T]$; Proposition~\ref{prop:sandwich49} gives the loss bound. A finite minimum of equally Lipschitz cones is equally Lipschitz. The identities $|z|=\operatorname{ReLU}(z)+\operatorname{ReLU}(-z)$ and $\min(u,v)=u-\operatorname{ReLU}(u-v)$ give the exact network. Finite lexicographic comparison and Borel repair give measurability. Effective refinement of the finitely many positively weighted terms yields the termination assertion.
\end{proof}

\subsection{Proof of the result \ref{thm:compile49}}\label{supp:proof-core49-3}
\begin{proof}
The triangle inequality bounds $f(x)$ above by each corner expression. Choose the lexicographically first attaining site $i_*(x)$. In every coordinate choose the cell endpoint between $x_j$ and $x_{i_*,j}$. At this corner distances add exactly. An owner of strictly smaller transformed value would improve the objective at $x$; a tied smaller owner would contradict the original tie rule. Thus this corner has owner $i_*(x)$, proving both identities. Along a coordinate line, forward and backward relaxations respectively minimize over sites on either side, always comparing value--original-index pairs. Iterating over coordinates gives the separable distance transform and its stated work. Rational arithmetic and unchanged repair complete the claim.
\end{proof}

\subsection{Proof of the result \ref{thm:regularity62}}\label{supp:proof-accuracy62-print-1}
\begin{proof}
The state cost has gradient infinity norm at most $8/d$, Hessian norm at most $22/d$, and diagonal second derivatives at most $5/d+16/d^2$. For the terminal cost the corresponding bounds are $10/d$, $26/d$, and $9/d+16/d^2$. They follow by differentiating the quadratic terms and the continuously differentiable squared shortage. The bound on the cycle Laplacian is four. The shortage Hessian, where it exists, is $16\mathbf1\mathbf1^\top/d^2$; the same second-difference bound holds across its kink.

For feasible $(x,a)$ and $(y,b)$, put $w=x-y$. The difference between the transition at their convex combination and the corresponding combination of transitions is $\theta(1-\theta)w_jw_{j+1}/16$ in coordinate $j$. Convexity and the $G_{t+1}$ Lipschitz bound therefore give a convexity defect of at most
\[
 \frac{\beta G_{t+1}}{16}\theta(1-\theta)
                 \sum_j|w_jw_{j+1}|
 \le\frac{\beta G_{t+1}}{16}\theta(1-\theta)\|w\|_2^2.
\]
The state cost supplies curvature $4/d$, while the two-control action cost has Hessian eigenvalues at least $7/4$. Thus the state curvature of $Q_t^*$ is at least $4/d-\beta G_{t+1}/8$, which is at least $107/(128d)$ when $G_{t+1}\le27/d$. Partial minimization over the affine capacity graph preserves the state convexity inequality. The minimizing action is unique, including at a binding capacity constraint.

To propagate Lipschitz continuity, write $a=c(x)u$, where $u$ lies in the fixed unit simplex. Transfer the same $u$ from $x$ to $y$. Each action-cost partial derivative has absolute value at most $13/16$, and $\sum_j u_j\le1$. The action-cost transfer contributes at most $13\|x-y\|_1/(128d)$. The transition Jacobian at a fixed action has column sums at most $11/16$; the capacity transfer adds at most $3/64$. Comparing a minimizing normalized action in both directions gives the stated recursion for $G_t$. Its right side maps $27/d$ below itself.

For a fixed $u$, the transition as a function of $x$ has Jacobian norm at most $11/16+1/16=3/4$. Its second-order composition term is bounded by $G_{t+1}/8$. The action-cost Hessian has norm at most $21/4$, and the capacity gradient has squared norm $1/(64d)$. These facts give the recursion for $M_t$. The infimum over $u$ preserves this common semiconcavity bound: after subtracting the same quadratic, each fixed-$u$ function is concave, and an infimum of concave functions is concave. The right side maps $54/d$ below itself. Convexity together with semiconcavity yields the claimed differentiability and gradient regularity; this conclusion does not require a differentiable optimizing-action map.

Along one state coordinate, the fixed-$u$ transition has zero second derivative. The squared norm of its first derivative is at most $85/256+23/(256d)$: the two original nonzero entries are bounded by $9/16$ and $1/8$, and the capacity-induced entries are bounded by $1/(16d)$. Combining this with the diagonal stage-cost bound gives $D_t$. Finally, positive interpolation along one coordinate has error between zero and $D_t h^2/8$. Applying the coordinate interpolations successively and adding the bounds proves \eqref{eq:interpolation62}.
\end{proof}

\subsection{Proof of the result \ref{thm:policy62}}\label{supp:proof-accuracy62-print-2}
\begin{proof}
Joint convexity, including the feasible graph, gives
\[
 Q_t^*(x,\bar\pi_t(x))\le\sum_vw_v(x)Q_t^*(v,a_t(v))
       \le I_nq_t^{\mathrm{sel},+}(x).
\]
The lower interpolation bound gives $V_t^*(x)\ge I_nl_t(x)-\kappa_t$. Subtraction proves the uniform optimal-advantage bound $e_t$. The partial action derivatives of $Q_t^*$ are bounded by $\Lambda_t$, so feasible implementation error adds at most $\Lambda_tb_t$. Adding the future policy loss and proceeding backwards proves the recursion. No convexity of $J^{\widehat\pi}$ is used.
\end{proof}

\subsection{Proof of the result \ref{lem:parabolic63}}\label{supp:proof-query63-1}
\begin{proof}
Apply concavity to $f(a)-Ha^2/2$ and substitute the two lower endpoints. The displayed quadratic has derivative $D-K+2Ks$. Minimizing it on the unit interval gives the three cases. Its coefficient of $K$ is $-s(1-s)$, which proves the last assertion.
\end{proof}

\subsection{Proof of the result \ref{thm:query63}}\label{supp:proof-query63-2}
\begin{proof}
At $r=1$, analytic endpoint enclosures and the minorant lemma give the assertion. Suppose it holds for $r-1$. Each child interval encloses the true optimal continuation. If the numerical next state is an interval, expanding the queried center value by $G_{r-1}$ times the maximum $\ell^1$ displacement encloses all its possible values. Conditional-mean lower bounds and the semiconcavity remainder then give valid endpoint brackets for the true $Q_r$. The minorants cover the entire action interval. A downward capacity endpoint can omit only an interval of length equal to the certified capacity-rounding difference; subtracting $\Lambda_r$ times that difference accounts for its possible improvement. The minimum upper endpoint is the cost of a feasible first action followed by the optimal future. This proves (i) by induction. Retaining earlier valid lower bounds preserves every conclusion.

For termination, an interval of length $h$ has minorant at least its smaller lower endpoint minus $H_rh^2/8$. Its smaller upper endpoint is an available global upper bound. With the endpoint-width guard, an interval of length at most
$h_\delta=\sqrt{3\delta/H_r}$ cannot be responsible for a gap exceeding $\delta$. Every refined parent is therefore longer than $h_\delta$. Its children are longer than $h_\delta/5$. Apart from the at most two initial fragments made by a prediction, disjoint terminal fragments consequently have total number at most $5c(x)/h_\delta$, up to endpoint counting. Equation~\eqref{eq:querycap63} is a conservative point budget. Bounded action-quantum rounding of a proposed quarter-to-three-quarter split preserves the one-fifth property whenever the quantum is smaller than $h_\delta/20$. Otherwise precision must be increased; a fixed-precision run is not allowed to ignore this requirement.

For (iii), let $\bar x$ be the acquired state. Its downward displacement is at most $d2^{-b}$ in $\ell^1$, and $c(\bar x)\le c(x)$ preserves feasibility. At a fixed action the stage cost is $8/d$-Lipschitz and the transition has column sums at most $11/16$. Thus replacing $\bar x$ by $x$ changes $Q_r$ by at most $(8+11\beta dG_{r-1}/16)2^{-b}$. The corresponding value change is at most $dG_r2^{-b}$. Add these two allowances to (i). Substituting the policy's own future yields the usual one-step loss recursion; backward induction gives \eqref{eq:querypolicy63}. No policy-value oracle and no population claim about fitting are used.
\end{proof}

\subsection{Proof of the result \ref{thm:gated64}}\label{supp:proof-gated64-1}
\begin{proof}
A routed action is merely another feasible partition point. The endpoint and minorant induction in Theorem~\ref{thm:query63} therefore applies without change. After each evaluation retain the smallest upper endpoint and the maximum of the old and newly valid global lower endpoints. Their difference cannot increase. The stopping test uses this difference, not a fitted value or a prediction error.

Put $h_\delta=\sqrt{3\delta/H_r}$. A partition interval of length at most $h_\delta$ has minimum minorant at least its smaller lower endpoint minus $3\delta/8$. The associated upper endpoint, with the stipulated width and capacity allowance, is within $\delta$ of that minorant. Such an interval cannot be responsible for an unresolved global gap. Every split parent therefore has length greater than $h_\delta$. The one-fifth condition makes each of its children longer than $h_\delta/5$. Starting with one interval, disjoint terminal fragments consequently number at most $1+5c(x)/h_\delta$. Counting endpoints and allowing for rounding of this numerical bound gives the displayed conservative budget. Floating failure is handled by the separately charged exact-rational procedure, not by this real-arithmetic count.

The first assertion about predictor evaluation follows from the order of the stopping test and the routing rule. For the last assertion, use backward induction on the remaining horizon, with the common analytic terminal kernel as base. At each nonterminal query, declined predictions leave the same conventional split, lower bound, witness and tie rule. Induction over partition refinements gives the same next query and returned endpoints; the induction over dates then gives the same implemented policy. Model-evaluation work can still differ, so transcript identity is not a zero-cost fitting claim.
\end{proof}

\subsection{Proof of Proposition \ref{prop:localization65}}\label{supp:proof-localization65-1}
\begin{proof}
Every retained lower endpoint is at least its true endpoint value minus $w$, and hence at least $f(a^*)-w$. On an interval of length at most $h$, the parabolic minorant is therefore at least $f(a^*)-w-Hh^2/8$. Subtracting the omitted-capacity allowance gives
$L\ge f(a^*)-w-Hh^2/8-\Lambda\Delta$; retaining an older, larger valid lower bound only improves this inequality. The upper endpoint at $\widehat a$ is at most $f(\widehat a)+w$. Semiconcavity and the one-sided first-order condition give
$f(\widehat a)-f(a^*)\le\lambda e+He^2/2$. Subtract the lower bound to prove \eqref{eq:localization65}. At an interior optimum the derivative is zero. The example $f(a)=a$ on $[0,c]$ shows that the boundary linear term is needed. Finally, retaining bounds makes $U'\le U$ and $L'\ge L$; expansion proves \eqref{eq:certificate-progress65}.
\end{proof}

\subsection{Proof of the result \ref{thm:conditional49}}\label{supp:proof-decisions49-1}
\begin{proof}
For each drawn bin sequence, conditional integration over its continuous realizations places the exact path expectation between the interval endpoints. Uniform mixing over the equal bins recovers the original continuous law. The lower endpoint expectation is thus no larger than $D_i$, and the upper endpoint expectation is no smaller. Apply the bounded empirical-Bernstein inequality of \citet{maurer2009} to the endpoint samples and take the union bound, with outward arithmetic included in the mean, variance and radius. Conditional validity for every $\mathcal H$ implies unconditional validity by integration. On this one simultaneous event, all ensuing deterministic comparisons of the covered quantities hold together.
\end{proof}

\subsection{Proof of the result \ref{cor:fee49}}\label{supp:proof-decisions49-2}
\begin{proof}
Subtract the gain intervals from $k$ and compare the resulting interval endpoints with zero. Taking the image of $[l,u]$ under the absolute-value map gives \eqref{eq:absolute49}. The assertions hold simultaneously for every nonnegative fee on the original coverage event.
\end{proof}

\subsection{Proof of the result \ref{thm:net49}}\label{supp:proof-decisions49-3}
\begin{proof}
Subtract the common $J_{b_0}$. The selected implementation has net difference at most $U(\lambda)$, and the best catalogue implementation has net difference at least $L(\lambda)$. This proves the regret bound pointwise on the one simultaneous event and hence for every $\lambda$. A minimum of finitely many affine functions changes its active affine expression only at an intersection. For $\lambda_2>\lambda_1$, sum the two minimizing inequalities to obtain $(\lambda_2-\lambda_1)(s_{\widehat b(\lambda_2)}-s_{\widehat b(\lambda_1)})\le0$. Since $s_b$ increases strictly with $b$, the precision comparative static follows.
\end{proof}
